% !TeX spellcheck = en_US
\documentclass[french]{yLectureNote}

\title{Mecanique Analytique}
\subtitle{Physique}
\author{Paulhenry Saux}
\date{\today}
\yLanguage{Français}
%lionels.calmels@cemes.fr
\professor{M.Brut}
\usepackage{graphicx}%----pour mettre des images
\usepackage[utf8]{inputenc}%---encodage
\usepackage{geometry}%---pour modifier les tailles et mettre a4paper
%\usepackage{awesomebox}%---pour les boites d'exercices, de pbq et de croquis ---d\'esactiv\'e pour les TP de PC
\usepackage{tikz}%---pour deiffner + d\'ependance de chemfig
% \usepackage{tabularx}%---pour dimensionner automatiquement les tableaux avec variable X
\usepackage{awesomebox}%---Pour les boites info, danger et autres
\usepackage{menukeys}%---Pour deiffner les touches de Calculatrice
\usepackage{fancyhdr}%---pour les en-t\^ete personnalis\'ees
\usepackage{blindtext}%---pour les liens
\usepackage{hyperref}%---pour les liens (\`a mettre en dernier)
\usepackage{caption}%---pour la francisation de la l\'egende table vers Tableau
\usepackage{pifont}
\usepackage{array}%---pour les tableaux
\usepackage{lipsum}
\usepackage{yFlatTable}
\usepackage{multicol}
\usepackage{tabularx}
\usepackage{booktabs}
\newcommand{\Lim}[1]{\lim\limits_{\substack{#1}}\:}
\renewcommand{\vec}{\overrightarrow}
\newcommand{\N}[0]{\mathbb{N}}
\newcommand{\dd}{\mathrm{d}}
\newcommand{\norm}[1]{||\vec{#1}||}
\newcommand{\fo}{\psi(\vec{r},t)}
\newcommand{\foe}{\psi(\vec{r},t)\*}
\newcommand{\HH}{\hat{H}}
\newcommand{\hb}{\hbar}
\newcommand{\lap}{\nabla^2}
\newcommand{\lapcc}{\frac{\partial^2 }{\partial x^2}+\frac{\partial^2 }{\partial y^2}+\frac{\partial^2 }{\partial z^2}}
\newcommand{\E}{\vec{E}(M)}
\newcommand{\B}{\vec{B}(M)}
\newcommand{\V}{V(M)}
\newcommand{\er}{\vec{e_{\rho}}}
\newcommand{\ep}{\vec{e_{\varphi}}}
\newcommand{\pip}{\pi^+}
\newcommand{\pim}{\pi^-}
\newcommand{\mv}{\mathcal{V}}
\DeclareMathOperator{\Div}{Div}
\newcommand{\rot}{\vec{\mathrm{rot}}}
\newcommand{\grad}{\vec{\mathrm{grad}}}
\newcommand{\qa}{q_{\alpha}}
\newcommand{\qdb}{\dot{q_{\beta}}}
\newcommand{\qb}{q_{\beta}}
\setcounter{chapter}{3}
\begin{document}
\chapter{Équation de hamilton}
\section{Notions de variable conjuguées}
On dédouble les coordonnées. On pose $Y_i = \dot{X_i}$.

On peut les lier aux transformations canoniques.

Question : Comment savoir si on peut arriver à un système d'équations connue. C'est légendre qui se la pose.

Il est utile de définir

% page 9 10 et 11 à faire
%--------------- APrès ces pages , séance d'aujourd'hui

On définit la transformée de Legendre : $g = -f+ux \Rightarrow \dd g = -v\dd y+x\dd y$
avec $-v = \frac{\partial g}{\partial y}, x = \frac{\partial g}{\partial u}$

On peut le comparer au lagrangien :
$$y = q_i, x = \dot{q_i}, u= p_i, f = L, g = H$$ avec $p_i = \frac{\partial L}{\partial \dot{q_i}}$

On a donc $$H(q_i, p_i, t) = \sum^N_{i=1}p_i\dot{q_i} - L$$

Puis $$\frac{\partial H}{\partial p_l}$$

On peut maintenant écrire :
$$H(q_i, p_l,t) = \sum_i \dot{q_i}p_i -L$$ et 2N équations 
$$\dot{q_l} = \frac{\partial H}{\partial p_l}$$ et $$\dot{p_j} = -\frac{\partial H}{\partial q_j}$$

On a des variables conjuguées.

\subsection{Exemple}
$$L = \frac12 m v^2 - V$$

$\frac{\partial L}{\partial \dot{r_i}} = \dot{mr_i} = p_i $

puis $H = \frac12 p_i \dot{r_i}+V = \sum_i \frac12 \frac{p_i^2}{m}+v$

On peut étudier dans l'espace des phases, où chaque point correspond à une condition initiale et 2 trajectoies ne peuvent pas se croiser.

% Truc sur l'extrmalisation de S : inutile

\subsection{Crochet de poisson}
Soient $q_\alpha$ et $p_{\beta}$ et 2 paires de fonctions $f$ et $g$
\begin{definition}
    $\{f,g\} = \sum_{\alpha} (\frac{\partial f}{\partial q_{\alpha}}\frac{\partial g}{\partial p_{\alpha}}-\frac{\partial f}{\partial p_{\alpha}}\frac{\partial g}{\partial q_{\alpha}})$
\end{definition}
Propriétés :
\begin{itemize}
    \item $\{f,g\} = - \{g,f\}$
    \item $\{f,gh\} = h\{f,g\} + g\{f,h\}$
    \item Identité de Jacobi : $\{f, \{g,h\}\}+ \{h, \{f,g\}\} + \{g, \{h,f\}\} = 0$
    \item lors de la trajectoire donnée par $H$ : $\frac{\dd }{\dd t}f(q_{\alpha}, p_{\alpha},t) = \{f, H\}+\frac{\partial f}{\partial t}$
\end{itemize}
$exercie avec L = n\times p$ 
\subsection{Théorèmes}
Soit un système dynamique à 2N variables $q_i, p_j$. la dynamique est donnée par $\dot{q_i} = F_i$ et $\dot{p_i} = G_i$

Existe-il un hamiltonien tq $F_i = \frac{\partial H}{\partial p_i}$ et $G_i =- \frac{\partial H}{\partial q_i}$

Si, $\forall f,g$, on a $\frac{\dd }{\dd t}\{f,g\} = \{\dot{f}, g\}+\{f,\dot{g}\}$

page 12

Théorème de Poisson :
Si fet g ne dépendent pas de t, sont des constantes du mouvement, alors $\{f,g\}$ esr aussi une constante du mouvement.

Corollaire : Si 2 composantes de L sont conservées, la troisème l'est aussi.

\subsection{Transformations canoniques}
Une transformation canonique préserve les crochets de poisson.

Notation symplectique

TODO 
Montrer que $\dot{\xi_i} = \sum_j  J_{ij}\frac{\partial H}{\partial \xi_j}$ et $\{f,g\} = \sum_{ij} \frac{\partial f}{\partial \xi_i}\frac{\partial g}{\partial \x_ih}J_{ij}$

page 13,14

Exercice du on pose : $Q_1 = q_1 \cos \alpha - \frac{\sin \alpha}{\beta}p_2$
$Q_2 = q_2 \cos \alpha - \frac{\sin \alpha}{\beta}p_1$
$P_1 = \beta q_2\sin \alpha +p_1\cos \alpha $, $P_2 = \beta q_1\sin \alpha +p_2\cos \alpha$
%TODO vérifeir que c'est une transformation canonique

Comme t trouver des transformations canoniques 

Remarque : si la transfo est indépendante du temps, \partial F/\partial t = 0 et H_2 = H()

\section{Théorèmes de Liouville}

\chapter{Méthodes de résolutions}
pq S = pas important


--- Partie Actions-angle pas aux programme



\end{document}
%CC = ex 10
